\scanpage{58}

\uline{Conclusions premières} (on pose $\overline{\mathbb R}=[-\infty,+\infty]$

\margin{\textcircled{1} Soit $X',X''\in\unsure{\mathcal E|\mathscr E}_\chi$ (resp. $\in\mathcal E_{\chi_!}$) $\Leftrightarrow X'\sqcup X''\in\mathcal E_\chi$ (resp. $\in\mathcal E_{\chi_!}$), \deleted{résulte de c et \unsure{e}} et alors $\chi(X)=\chi(X')+\chi(X'')$ (resp. $\ldots$) résulte de c) et e) ($\varnothing\in\mathcal E_\chi$) resp. de b)}

\textcircled{\unsure{1|2}} \deleted{\unsure{Soit|Pour} \unsure{$I$}} \inserted{$\forall n\in\mathbb N$}, $R^n,\overline R^n,B^n,S^n\in\mathcal E_\chi\cap\mathcal E_{\chi_!}$, de plus
\[
\left\{\begin{aligned}
\chi(\mathbb R^n)=\chi(\overline R^n)&=\chi(B^n)=1 & \chi(S^n)&=\chi_!(S^n)=1+(-1)^n\\
\overset{\Vert}{\chi_!}(\overline R^n)&=\overset{\Vert}{\chi_!}(B^n)\\
\chi_!(R^n)&=(-1)^n
\end{aligned}\right.
\]

\uline{Dém} $R^n,\overline R^n,B^n\sim\mathrm{pt}$, donc par d) et e) ils sont $\in\mathcal E_\chi$ ; comme $\overline R^n,B^n$ compacts, par b) ils sont aussi $\in\mathcal E_{\chi_!}$.

On a $S^n=S^n_+\cup S^n_-$ et $S^n_+\cap S^n_-\simeq S^{n-1}$, \inserted{Or $S^n_+\simeq S^n_-\simeq B^n$ donc $\in\mathcal E_\chi$} donc on conclut de c) par réc. sur $n$ que $S^n\in\mathcal E_\chi$ (\uline{NB} $S^0=\{\mathrm{pt}\}\sqcup\{\mathrm{pt}\}\in\mathcal E_\chi$ par d) et e), et $\chi(S^0)=1+1=2$), et on trouve $\chi(S^n)=\chi(S^n_+)+\chi(S^n_-)-\chi(S^{n-1})\deleted{\unsure{+}}=2-\chi(S^{n-1})$, d'où par récurrence sur $n$ $\chi(S^n)=1+(-1)^n$. Comme $S^n$ cpt, $S^n\in\mathcal E_{\chi_!}$ et $\chi_!(S^n)=\chi(S^n)$ par (a). Enfin, $R^n\simeq S^n-\mathrm{pt}$, donc par (b) $R^n\in\mathcal E_{\chi_!}$ et $\chi(R^n)=\chi(S^n)-\chi(\mathrm{pt})=(1+(-1)^n)-1=(-1)^n$, cqfd.

\textcircled{3} Soit $X_{-1}=\varnothing\subset X_0\subset X_1\subset\cdots\subset X_n=X$ filtration de l'espace $X$, tel que $\forall 0\leq i\leq n$, $X_i-X_{i-1}$ \deleted{\unsure{est} \unsure{$\mathbb R$}} a un nb fini $f_i$ de comp. connexes ttes homéom. à $\mathbb R^i$ (i.e. $X_i-X_{i-1}\simeq f_i\times\mathbb R^i$). Alors $X\in\mathcal E_{\chi_!}$, et on a :
\[
\chi_!(X)=\sum(-1)^i f_i
\]
(donc si $X$ est cpt, on a $\deleted{\unsure{\chi}}X\in\mathcal E_\chi$ et $\chi(X)=\sum(-1)^i f_i\ \ldots$

% archive sheet number: 57

